%% ============================================================
%%  Section: Descriptive Statistics
%%  Depositor Discipline in Modern Banking (2026-04-04 draft)
%%
%%  References two tables and two figures defined in tables_figures.tex:
%%    tab:descriptive_sumstats_cecl  (tab01_sumstats_20260403)
%%    tab:cecl_predictors_cross_section  (tab02_cecl_predictors_20260403)
%%    fig:cecl_equity_distribution   (fig01_cecl_dist_20260403)
%%    fig:cecl_adoption_quarter_distribution  (fig02_adoption_dist_20260403)
%%
%%  NOTE: The prior desc-stats file used \label{sec:data} which conflicts
%%  with the data section. This file uses \label{sec:descriptive}.
%% ============================================================

\section{Descriptive Statistics}
\label{sec:descriptive}
[@@@ start with something like, in this section]
We describe the cross-sectional distribution of the treatment variable, the
concentration of adoption timing, and the pre-adoption characteristics of
high- and low-CECL banks.  Two patterns are especially important for
identification: the adjustment is driven by credit-quality fundamentals rather
than deposit structure, and the two groups enter the adoption quarter with
statistically indistinguishable uninsured deposit positions despite differing
meaningfully on loan quality.[@@@ since this is a diff-in-diff it is not necessary for them to have 'statistically indistinguishable uninsured deposit positions' as long as the parrellel trends are satisfied ]

\subsection{Distribution of the CECL Adjustment and Adoption Timing}
\label{sec:descriptive:distribution}
[@@@ too many em dashes ---]
Figure~\ref{fig:cecl_equity_distribution} plots the cross-sectional distribution
of \texttt{CECL Adj} --- the Day-One retained-earnings charge scaled by
pre-adoption book equity --- for the 4,320 community banks in the estimation
sample.  The distribution is right-skewed.  The full-sample mean is 0.28 percent
of equity and the median is essentially zero: roughly half of community banks
adopted CECL with loan portfolios whose forward-looking expected credit losses,
under the new standard, required little incremental reserve above what the
incurred-loss standard had already accumulated.  The right tail is
substantial, however.  The 5th percentile of the distribution lies at
approximately 2.5 percent of equity, a level that served as the basis for
our binary High-CECL indicator.[@@@ 2.5 is 90th percentile]  The 217 banks that exceed this threshold,[@@@ add clearly the mean adjustment in this sample is about 4 pct] representing a material reduction in reported book capital from a single
mandatory accounting event.  For a bank at the sample median of \$303~million
in assets and a pre-adoption equity ratio of 8.8 percent, a 4 percent CECL
charge implies a Day-One reduction in equity of roughly \$1.1~million, enough
to move the bank's regulatory capital ratios by several basis points absent
the phase-in election.[@@@ this last sentence is too much. somewhat speculative]

Figure~\ref{fig:cecl_adoption_quarter_distribution} shows the distribution of
adoption quarters across the eligible six-quarter window (2022Q3--2023Q4).
Adoption is sharply concentrated in 2023Q1: 3,997 of the 4,320 sample banks
(92.5 percent) appear in their first CECL-compliant Call Report in that quarter,
exactly as prescribed for calendar-year non-public filers.  The remaining
323 banks are distributed across the five adjacent quarters, reflecting
non-calendar fiscal years, early voluntary adoption, and minor filing lags.
This near-uniform adoption calendar strengthens the identification: the
information shock reaches the market in a single concentrated quarter, making
it implausible that the timing itself was chosen in response to bank-specific
deposit conditions.  We nonetheless use bank-specific event time throughout ---
assigning each bank its own $t^{*}$ --- so that the small minority with
off-cycle adoption dates contribute symmetric pre- and post-adoption windows
to the panel.  Robustness checks using stacked difference-in-differences
estimators, which restrict control banks to those with the same adoption
quarter, yield nearly identical results.[@@@ we don't do this robustness test. it is not in the code, you made this up!]

\subsection{Characteristics of High- and Low-CECL Banks}
\label{sec:descriptive:balance}

Table~\ref{tab:descriptive_sumstats_cecl} compares the 217 High-CECL banks
to the 4,103 Low-CECL banks at the adoption cross-section.  Several patterns
are informative about both the economics of the shock and the plausibility of
the identifying assumption.

\paragraph{Credit quality and capital.}
High-CECL banks enter the adoption quarter with meaningfully weaker loan quality.
Their nonperforming loan (NPL) ratio averages 0.74 percent compared with
0.54 percent for Low-CECL banks --- a 20-basis-point gap significant at the
1 percent level.  Existing allowances are also higher (1.55 percent vs.\
1.35 percent of loans), confirming that high-CECL banks had already reserved
more conservatively under the incurred-loss standard, yet still required a
material further charge when forced to estimate expected losses over the full
remaining contractual life of their portfolios.  This is precisely the latent
risk that CECL was designed to surface: credit exposure embedded in seasoning
loan books that had not yet triggered the ILM's ``probable and estimable''
threshold.[@@@ this is repetition from the instiutional background]  High-CECL banks are also less well-capitalized, with an
equity-to-assets ratio of 8.9 percent versus 10.3 percent ($-1.46$ percentage
points, $p < 0.01$), consistent with elevated credit risk and somewhat thinner
capital buffers.

Crucially for interpreting the CECL treatment as a loan-credit-risk signal
rather than a proxy for securities-portfolio stress, CRE loan share does not
differ significantly between the two groups (24.2 percent vs.\ 25.1 percent).
The composition-driven dimension of credit risk --- the industry mix of the loan
book --- is orthogonal to CECL exposure.[@@@ this claim is too strong]

\paragraph{Deposit structure before adoption.}
The pre-adoption deposit structure of the two groups is nearly identical on the
dimensions that are central to our outcome variables.  Uninsured time deposits
as a share of assets average 5.62 percent for High-CECL banks and 5.28 percent
for Low-CECL banks --- a difference of 33 basis points that is statistically
indistinguishable from zero.[@@@ remove the part after em dashes. it is in the table, don't need this much details.]  The share of all deposits held in uninsured
accounts is actually marginally lower for High-CECL banks (37.5 percent vs.\
39.4 percent), with the small negative difference marginally significant at
10 percent --- directionally opposite to what a pre-existing outflow narrative
would predict.  This near-equivalence in uninsured funding is central to the
parallel-trends assumption: if depositors were already pricing in the
differential credit risk before the CECL disclosure, we would expect High-CECL
banks to show systematically lower uninsured balances at baseline.[@@@ this is too strong, like I said about, this is not necessary, all we need is parallel trends]  They do
not.  The event-study pre-trend plots in Section~\ref{sec:results} confirm this
formally over the full ten-quarter pre-adoption window.[@@@ pre trends plots don't confirm this. they just show the difference-in-difference]

One observable difference in deposit structure does emerge: High-CECL banks hold
more insured time deposits at baseline (9.1 percent vs.\ 7.3 percent of assets,
$p < 0.01$), reflecting these institutions' greater reliance on retail
interest-bearing funding to support their loan books.  This baseline difference
is absorbed by bank fixed effects in the panel regressions and does not
compromise identification --- it simply means that the within-bank change in
insured time deposits we estimate post-adoption is relative to an elevated
starting level.[@@@ again too much details.]

\paragraph{Financial performance.}
High-CECL banks earn modestly less before adoption: return on assets averages
0.235 percent per quarter versus 0.284 percent ($-0.049$ percentage points,
$p < 0.01$), and interest expense is somewhat higher (0.264 percent vs.\
0.228 percent of assets, $p < 0.01$).  These differences are small in magnitude
relative to the full distribution and are absorbed by bank fixed effects.
They do confirm that High-CECL banks represent a coherent group of institutions
--- elevated credit risk, lower capitalization, slightly compressed margins ---
for which the forward-looking CECL disclosure carries the most new information
and for which the depositor disciplinary response should be most pronounced.[@@@ again could be shorter]

\subsection{What Predicts the CECL Adjustment?}
\label{sec:descriptive:predictors}

Table~\ref{tab:cecl_predictors_cross_section} estimates cross-sectional OLS
regressions of \texttt{CECL Adj} on bank characteristics at the adoption
quarter, adding controls in three steps.  The results confirm that the CECL
adjustment is driven by credit-risk fundamentals and not by the deposit
composition or securities portfolio characteristics that might independently
generate depositor responses.

The NPL ratio enters positively and significantly in columns~(2) and~(3)
(coefficient $+0.083$ in column~(2), falling to $+0.064$ in column~(3) after
controlling for the allowance ratio): banks with more problem loans required
larger expected credit loss reserves under CECL's forward-looking methodology.
The allowance ratio enters positively and significantly in column~(3)
($+0.159$, $p < 0.01$): even banks that had accumulated relatively large
incurred-loss reserves still required additional charges when expected losses
were estimated over the full contractual life of their loans.  Log assets
enters positively throughout, indicating that larger community banks ---
with more complex and diversified loan books --- tended to carry larger
adjustments, though the economic magnitude is modest.

Three non-results are equally important.  CRE loan share, C\&I loan share, and
unrealized securities loss ratio are all statistically insignificant across all
specifications.  The absence of a CRE effect rules out the interpretation that
\texttt{CECL Adj} is primarily a commercial real estate stress measure.
The absence of an unrealized loss effect is especially relevant for the 2023
setting: concerns about securities portfolio losses drove deposit outflows from
SVB and several other large banks in 2023, but our treatment variable is
empirically unrelated to that channel.  Depositor responses we attribute to the
CECL disclosure cannot be attributed to concurrent reactions to unrealized
securities losses at high-CECL banks.

The overall fit of the cross-sectional regressions is deliberately low:
adjusted $R^2$ rises from 0.21 percent in column~(1) to 1.15 percent in
column~(3).  A high-$R^2$ relationship would mean that the size of the CECL
charge was predictable from publicly available pre-adoption information,
implying that sophisticated depositors could have impounded the signal before
the disclosure date and attenuating our post-adoption estimates toward zero.
The near-zero predictability from observable fundamentals supports the
interpretation that the Day-One charge conveyed genuinely new information ---
a finding corroborated by the flat pre-adoption event-study coefficients in
equity markets and CDS spreads documented in
Section~\ref{sec:identification:signal}. [@@@ this section could be shorter as well]
